\documentclass[10pt,a4paper]{amsart}
\usepackage[margin=19mm]{geometry}
\usepackage{amsmath,amssymb,mathtools}
\usepackage[scr=rsfs,cal=euler]{mathalfa}
\usepackage{xfrac}
\usepackage[unicode,hidelinks,bookmarks=false]{hyperref}
\newcommand{\ie}{i.e.\ }
\newcommand{\cf}{cf.\ }
\newcommand{\resp}{respectively\ }
\newcommand{\Subsection}{\subsection}
\providecommand{\nobreakdash}{\nobreak\hbox{-}}
\setlength{\emergencystretch}{2em}
\multlinegap0pt
\usepackage{mathtools}
\usepackage{amssymb}
\usepackage{xcolor}
\usepackage{tikz}
\usepackage{enumerate}
\usepackage[shortlabels]{enumitem}
\usepackage{dsfont}
\usepackage{float}
\usepackage{bm}
\usepackage{algorithm}\usepackage{algpseudocode}
\newtheorem{lemma}{Lemma}
\newtheorem{claim}{Claim}
\usepackage{cleveref}
\crefname{lemma}{Lemma}{Lemmas}
\crefname{claim}{Claim}{Claims}
\title{\bf\Large Tur\'{a}n numbers of $4$-uniform tight even cycles minus one edge}
\author{Wanfang Chen, Jianfeng Hou, Xizhi Liu, Yixiao Zhang, Hongbin Zhao}
\date{Source: arXiv:2606.22828v1 (CC BY 4.0)}
\begin{document}
\maketitle

\noindent\emph{Source and licence.} \bf\Large Tur\'{a}n numbers of $4$-uniform tight even cycles minus one edge. Version arXiv:2606.22828v1 (CC BY 4.0), distributed by arXiv under the Creative Commons Attribution 4.0 International licence, http://creativecommons.org/licenses/by/4.0/, as recorded in the arXiv licence metadata for this exact version and shown on the abstract page https://arxiv.org/abs/2606.22828v1. Full licence notice: \texttt{licenses/arxiv-2606.22828-CC-BY-4.0.txt}.\par\smallskip

\noindent\textit{Editorial scope.} Counted unit: the article's lemma lem:J-free-local-link-stability (source0.tex lines 999--1003). The article's complete original proof of it is delivered with it (source0.tex lines 1005--1091, one \texttt{proof} environment of 87 lines). No mathematics is written, repaired or re-derived: the statement and the proof are the article's own lines, in the article's own notation, and no retained line is substituted or re-flowed.  No further source-local context is needed: every cross-reference of every retained line resolves inside the two delivered spans.  Nothing was omitted from any retained span, so the package carries no omission record.  No retained line contains a citation, so the wrapper carries no bibliography and no bibliography entry.  No retained line contains a figure environment, an image-inclusion command or drawing code, so no image is delivered and the package declares no asset dependency.  Editorial formatting only, in addition: an AMS \texttt{amsart} preamble that restates the article's own declarations and macros used by the retained lines, namely the environments \texttt{lemma}, \texttt{claim} on separate counters, together with the standard packages those lines need.  The retained environments print in this document as Lemma 1, Claim 1 in order of appearance, whereas the article prints the same items with the numbering of its own full document.  The complete native source of the article as distributed in the arXiv e-print for this exact version is bundled unchanged as \texttt{sources/203385/source0.tex}.
\begin{lemma}\label{lem:J-free-local-link-stability}
For every $\xi_0>0$, there exist $\beta>0$ and $n_{\mathcal J}$ such that the following holds.
Let $L$ be a $\mathcal J$-free two-colored $3$-graph on $n\ge n_{\mathcal J}$ vertices with color classes $A\sqcup B$.
If $\bigl||A|-|B|\bigr|\le \beta n$ and $|L|\ge \left(\frac12-\beta\right)\binom n3$, then $\min\Bigl\{|L\triangle \mathsf O(A,B)|,\ |L\triangle \mathsf E(A,B)|\Bigr\} \le \xi_0 n^3$. 
\end{lemma}

\begin{proof}
Let $V\coloneqq V(L)$.
Let $0<\beta\ll\eta\ll\xi_0$, and suppose that $L$ is a $\mathcal J$-free two-colored $3$-graph with color classes $A\sqcup B$, satisfying $\bigl||A|-|B|\bigr|\le \beta n$ and $|L|\ge \left(\frac12-\beta\right)\binom n3$.
For a pair $e$, write $N_A(e)\coloneqq N_L(e)\cap A$ and $N_B(e)\coloneqq N_L(e)\cap B$. Since $L$ is $\mathcal J$-free, every pair has its link contained in one color class: $N_A(e)=\varnothing$ or $N_B(e)=\varnothing$.
Indeed, for an $AA$-pair this is exactly the exclusion of $J_1$, for a $BB$-pair it is the exclusion of $J_1^{\mathrm{sw}}$, and for a mixed pair it follows from the exclusions of $J_2$ and $J_2^{\mathrm{sw}}$.

Call a pair $e$ \emph{good} if $d_L(e)\ge (1/2-\eta)n$.

\begin{claim}\label{clm:many-good-pairs}
All but $2\beta\eta^{-1}n^2$ pairs are good.
\end{claim}

\begin{proof}
Since every pair has degree at most $\max\{|A|, |B|\} \le (1/2+\beta)n$ (by the previous paragraph), and taking $n_{\mathcal J}\ge\beta^{-1}$, double counting gives
\[
\begin{aligned}
\sum_{e\in\binom V2}\left(\left(\frac{1}{2}+\beta\right)n-d_L(e)\right)
=\left(\frac{1}{2}+\beta\right)n\binom n2-3|L| 
\le \left(\frac{1}{2}+\beta\right)n\binom n2-3\left(\frac{1}{2}-\beta\right)\binom n3 
\le 2\beta n^3.
\end{aligned}
\]
If $M$ is the number of non-good pairs, then each non-good pair $e$ contributes at least $\eta n$ to the last sum, because $d_L(e)<(\frac{1}{2}-\eta)n$.
Thus $M\eta n\le 2\beta n^3$, so $M\le 2\beta\eta^{-1}n^2$.
\end{proof}

For a good pair $e$, let $\tau(e)\in\{A,B\}$ be the unique color class containing its link.
Then $N_L(e)\subseteq \tau(e)$ and, by goodness, $|N_L(e)|=d_L(e)\ge(\frac{1}{2}-\eta)n$.
On the other hand, the balance assumption gives $|\tau(e)|\le(\frac{1}{2}+\beta)n$.
Therefore
\[
|\tau(e)\setminus N_L(e)| \le |\tau(e)|-|N_L(e)| \le (\beta+\eta)n.
\]

\begin{claim}\label{clm:label-propagation}
Let $xy$ be a good same-color pair.  Every good mixed pair sharing $x$ or $y$ has label different from $\tau(xy)$.
\end{claim}

\begin{proof}
By symmetry, suppose $x,y\in A$.  If $\tau(xy)=B$ and some good mixed pair $xb$ has label $B$, then
\[
|N_L(xy)\cap B|+|N_L(xb)\cap B|
\ge 2\left(\frac{1}{2}-\eta\right)n>|B|
\]
for $\eta,\beta$ sufficiently small.  Hence there is $c\in B$ such that $xyc,xbc\in L$, which gives a copy of $J_2$, a contradiction.  Thus every good mixed pair $xb$ has label $A$.  The same argument applies to $yb$.

If $\tau(xy)=A$ and some good mixed pair $xb$ has label $A$, then
\[
|N_L(xy)\cap A|+|N_L(xb)\cap A|
\ge 2\left(\frac{1}{2}-\eta\right)n>|A|,
\]
so there is $a\in A$ such that $xya,xba\in L$, giving a copy of $J_1$.  This is again impossible.  The proof with $A$ and $B$ interchanged is identical.
\end{proof}

\begin{claim}\label{clm:same-color-common-label}
There are subsets $A_0\subseteq A$ and $B_0\subseteq B$ and labels $s_A,s_B\in\{A,B\}$ such that
\[
|A\setminus A_0|+|B\setminus B_0|=O(\beta\eta^{-2}n+\eta n),
\]
every good pair in $\binom{A_0}{2}$ has label $s_A$, and every good pair in $\binom{B_0}{2}$ has label $s_B$.
\end{claim}

\begin{proof}
We prove the assertion for $A$; the proof for $B$ is the same.  Let $A'$ be the set of vertices $x\in A$ incident with at most $\eta n$ non-good pairs inside $A$.  By \cref{clm:many-good-pairs}, we have $|A\setminus A'|\eta n\le 2\cdot 2\beta\eta^{-1}n^2$, and hence $|A\setminus A'|=O(\beta\eta^{-2}n)$.
For every $x\in A'$, the vertex $x$ is incident with a good pair inside $A$.

If $x\in A'$ is incident with good $AA$-pairs of both labels, then \cref{clm:label-propagation} gives incompatible labels for every good mixed pair from $x$ to $B$.  Hence every pair $xb$ with $b\in B$ is non-good.  There are at most $2\beta\eta^{-1}n^2$ non-good pairs altogether, so the number of such vertices is $O(\beta\eta^{-1}n)$, which is $O(\beta\eta^{-2}n)$.  After deleting them from $A'$, every remaining vertex $x$ has a well-defined label $\lambda_A(x)$, namely the common label of the good $AA$-pairs incident with $x$.

Let $A_A$ and $A_B$ be the remaining vertices with $\lambda_A(x)=A$ and $\lambda_A(x)=B$, respectively.  If both sets had size at least $\eta n$, then every pair crossing $A_A$ and $A_B$ would be non-good, since a good crossing pair would have to have both labels.  This would give at least $\eta^2n^2$ non-good pairs, contradicting \cref{clm:many-good-pairs} because $\beta\ll\eta^3$.  Thus one of $A_A,A_B$ has size less than $\eta n$.  Delete this smaller set and call the remaining set $A_0$; its common label is $s_A$.
\end{proof}

\begin{claim}\label{clm:mixed-labels-and-template}
There is a label $s\in\{A,B\}$ such that all but $O(\eta n^2)$ pairs have the following labels: same-color pairs have label $s$, and mixed pairs have the other label.  Consequently, $\min\Bigl\{|L\triangle \mathsf O(A,B)|,\ |L\triangle \mathsf E(A,B)|\Bigr\}=O(\eta n^3)$. 
\end{claim}

\begin{proof}
Let $A_0,B_0,s_A,s_B$ be as in \cref{clm:same-color-common-label}.  All pairs meeting $(A\setminus A_0)\cup(B\setminus B_0)$, together with all non-good pairs, account for only $O(\eta n^2)$ pairs, using $\beta\ll\eta^3$.  Since $|A_0||B_0|\ge n^2/5$ for $n$ large and $\eta$ small, there is at least one good mixed pair between $A_0$ and $B_0$.

Consider a good mixed pair $ab$ with $a\in A_0$ and $b\in B_0$.  Since $a$ is incident with a good $AA$-pair of label $s_A$, \cref{clm:label-propagation} forces $\tau(ab)\ne s_A$.  Similarly, since $b$ is incident with a good $BB$-pair of label $s_B$, it forces $\tau(ab)\ne s_B$.  Applying this to one good mixed pair between $A_0$ and $B_0$ gives $s_A=s_B$; call this common label $s$.  It follows that all good mixed pairs between $A_0$ and $B_0$ have the label different from $s$.

Thus all but $O(\eta n^2)$ pairs have the asserted labels.  If $s=B$, let $\Pi_s=\mathsf O(A,B)$, and if $s=A$, let $\Pi_s=\mathsf E(A,B)$.  Every edge of $L\setminus\Pi_s$ contains a pair which is either non-good or has the wrong label, so $|L\setminus\Pi_s|=O(\eta n^3)$.

Conversely, apart from triples containing a non-good or wrongly labeled pair, a triple of $\Pi_s$ can be missing only because one of its correctly labeled good pairs misses the third vertex.  For each good pair $e$, the number of such missed vertices is at most $(\beta+\eta)n=O(\eta n)$.  Summing over all pairs gives $|\Pi_s\setminus L|=O(\eta n^3)$.
\end{proof}

Choosing $\eta\ll\xi_0$ and then $\beta\ll\eta^3$, and taking $n$ sufficiently large, \cref{clm:mixed-labels-and-template} proves the lemma.
\end{proof}

\end{document}
